\section{引言}
\label{sec:introduction}

在二维光子晶体中引入一条线缺陷，会破坏垂直于缺陷方向的周期性，但保留沿缺陷方向的
周期性。当频率落在背景晶体的谱隙中时，这样的缺陷可能支持沿缺陷传播、而在横向指数
局域的电磁波，即\emph{导模}\cite{JoannopoulosJohnsonWinnMeade2008}。固定沿缺陷方向的
Bloch 波数 $\beta$ 后，导模对应于一个无界条带上自伴算子 $A(\beta)$ 位于本质谱隙中的
特征值\cite{Fliss2013}。这既不是完整周期介质的能带计算，也不是均匀开放波导的截面
特征值问题：特征值是由缺陷诱导、嵌在无界周期背景中的离散谱。

\paragraph{计算上的两个困难.}
第一个困难是横向无界。超胞方法把条带截断为有限宽度并施加周期条件，其模型误差随导模
局域程度变差而迅速增大\cite{Fliss2013,KlindworthSchmidtFliss2014}。Joly、Li 与 Fliss
用单胞问题构造半无限周期波导的精确边界条件\cite{JolyLiFliss2006}；Fliss 据此把导模
问题等价地约化到缺陷附近的有界区域，代价是得到一个关于频率非线性的特征值问题
\cite{Fliss2013}。这一框架已有高阶有限元实现、Robin-to-Robin 变体和色散曲线追踪
\cite{KlindworthSchmidtFliss2014,FlissKlindworthSchmidt2015,Klindworth2015}；基于单胞
DtN 映射的约化也被用于光子晶体波导和缺陷模的计算
\cite{HuangLuLi2007,LiLu2007,YuanLu2007}。

第二个困难是可靠性。对分片常数介质，边界积分方法只需离散材料界面，并可借助准周期
Green 函数自动满足纵向准周期条件，因而能以很少的未知量达到高阶精度
\cite{Linton1998,BarnettGreengard2010,Kress1991,HaoBarnettMartinssonYoung2014}。但在
非线性特征值问题中，离散矩阵近乎奇异只说明某个有限维方程近似可解
\cite{GuettelTisseur2017}；它并不直接给出连续算子特征值的误差，更不能证明连续谱点的
存在。有限元框架中已有针对光子晶体特征值的后验估计
\cite{Giani2013,EngstromGianiGrubisic2016}；Barnett、Hassell 与 Tacy 则表明，对有界
光滑区域上的 Neumann 问题，试验函数的边界不相容量可以直接控制到最近谱点的距离
\cite{BarnettHassellTacy2018}。本文的问题横向无界、含尖锐材料界面，且离散场在半波导中
由传播关系逐胞元生成，这些结果都不能直接套用。

\paragraph{本文的工作.}
本文研究两侧为周期介质（左右可以不同）、中心为局部缺陷的二维条带波导。介质采用分片
常数系数，每个周期胞元恰含一个单连通介质柱，中心胞元可以含一个形状不同的介质柱，也可以为空。主要
内容如下。
\begin{enumerate}
  \item \emph{以 Cauchy 数据为主线的中心约化}（第~\ref{sec:reduction}~节）。把半波导
    的全部有限能量解在人工界面上的 Cauchy 数据收集为一个子空间，证明导模等价于中心胞元
    中 Cauchy 数据落在两个子空间中的非零解。这一表述不要求半波导 Dirichlet 问题唯一
    可解，DtN 形式只是其子空间为图像时的特例。
  \item \emph{边界积分与端口 Rayleigh 展开的离散}（第~\ref{sec:discretization}~节）。
    周期胞元用准周期层势和传输边界积分方程离散，并在端口用 Rayleigh 振幅表示，得到
    单胞散射矩阵；有序 QZ 分解给出衰减子空间与胞元传播矩阵。中心胞元的场由层势与均匀
    背景 Rayleigh 波表示，最终得到一个以 Cauchy 数据匹配形式写出的离散非线性特征值
    问题。消去中心未知量后得到端口约化矩阵 $\mathsf A(k)$；它除极点外解析，候选
    $\widehat k$ 由 $\det\mathsf A(k)$ 的 Chebyshev 插值求根得到\cite{Boyd2002,ZhaoBarnett2015}。
  \item \emph{谱距离定理}（第~\ref{sec:error}~节）。用候选系数和连续核在整个条带上构造
    一个分片场，它逐片精确满足 Helmholtz 方程，只在材料界面和人工竖直界面上留有跳跃。
    把界面按几何分为若干组，对每组用领子上的迹估计和 Green 恒等式，证明
    $\operatorname{dist}(\widehat k^2,\sigma(A(\beta)))$ 受各组跳跃的加权 $L^2$ 范数
    $\varepsilon_g$ 与迹常数 $C_g$ 的乘积之和控制（附录~\ref{app:trace}）。这一结论关于整个谱
    成立，不需要预先验证本质谱隙。
\begin{workonly}
原稿第 3 项：“利用正移位预解式的全界面迹估计和 Green 恒等式，证明
$\operatorname{dist}(\widehat k^2,\sigma(A(\beta)))$ 受这些跳跃的加权 $L^2$ 范数控制。”
\end{workonly}
  \item \emph{计入舍入的可靠计算}（第~\ref{subsec:computation}~小节、附录~\ref{app:reliable}）。
    跳跃积分按面板用 Chebyshev 插值误差界和球算术计算，无穷多个周期胞元的贡献由传播
    矩阵的收缩性求和；场范数的下界取自若干圆盘上的质量，迹常数由领子几何条件的逐段
    验证给出。全部输入都计入舍入，因此所得区间是可靠的谱距离上界。
  \item \emph{带隙证书}（第~\ref{subsec:gap-certificate}~小节）。用有限个平面波矩阵的惯性
    测试和 Bloch 能带关于相位的 Lipschitz 界，证明两侧引线的公共谱隙。区间落在其中时，
    其中至少有一个导模特征值。
  \item \emph{数值例子}（第~\ref{sec:results}~节）。对中心为星形柱、两侧为不同三叶柱
    引线的例子，一次完整运行在 $14$ 分钟内给出全宽 $2.74\times10^{-11}$ 的可靠波数区间，
    其中至少有一个导模特征值；区间约为同一问题 FEM 区间宽度的 $1/163$。
\begin{workonly}
原稿第 4 项：“\emph{可计算的界}。说明如何计算材料界面与人工界面上的跳跃积分，如何用传播
矩阵的收缩性控制无穷多个周期胞元的贡献，以及如何把各输入替换为单侧可靠界。”
\end{workonly}
\end{enumerate}

\paragraph{本文不声称的内容.}
\pending{2026-11-06：带隙证书已加入，改写本段}
谱距离界只说明以 $\widehat k^2$ 为中心、给定半径的区间与 $\sigma(A(\beta))$ 相交；加上
带隙证书，可知区间中至少有一个孤立、有限重的特征值。本文不声称该特征值在区间中唯一，
也不确定它的编号或重数。离散方法的先验收敛性也不在本文讨论范围之内：定理只作用于一个
给定的候选和重构场，与候选是如何得到的无关。

\begin{workonly}
原稿：
谱距离界只说明以 $\widehat k^2$ 为中心、给定半径的区间与 $\sigma(A(\beta))$ 相交。
只有另行证明该区间落在本质谱隙内，才能把其中的谱点解释为离散导模特征值；本文不据此
声称特征值唯一、编号已知或谱隙本身已被证明。离散方法的先验收敛性也不在本文讨论范围
之内：定理只作用于一个给定的候选和重构场，与候选是如何得到的无关。
\end{workonly}

\paragraph{结构.}
第~\ref{sec:modelproblem}~节给出模型和算子表述；第~\ref{sec:reduction}~节给出连续的
中心约化；第~\ref{sec:discretization}~节给出离散与候选计算；第~\ref{sec:error}~节构造
重构场并证明谱距离界；第~\ref{sec:results}~节给出数值结果；附录~\ref{app:center-cell}
证明中心胞元表示定理，附录~\ref{app:trace} 给出分组迹估计的几何与证明，
附录~\ref{app:reliable} 给出计入舍入的可靠计算。

\paragraph{记号.}
表~\ref{tab:notation} 列出全文的主要记号。复数内积对第二个变量取共轭；
$\lVert\cdot\rVert$ 无下标时指 $\mathcal H$ 范数或矩阵 $2$-范数，由上下文确定。

\begin{longtable}{p{0.27\textwidth}p{0.66\textwidth}}
  \caption{主要记号}\label{tab:notation}\\
  \toprule
  记号 & 含义\\
  \midrule
  \endfirsthead
  \toprule
  记号 & 含义\\
  \midrule
  \endhead
  \bottomrule
  \endfoot
  $\beta$, $d$, $L_\pm$ & $y$ 方向 Bloch 波数与周期；左右引线的 $x$ 方向周期\\
  $k$, $\mu=k^2$ & 波数与谱参数；候选记为 $\widehat k$\\
  $\rho$, $\rho_{\mathrm e}$, $\rin,\rax$ & 分片常数系数，背景值，正的上下界\\
  $B$, $I_d$ & 条带 $\RR\times I_d$，$I_d=(-d/2,d/2)$\\
  $A(\beta)$, $\mathcal H$ & 条带上的自伴算子及加权空间 $L^2(B,\rho\,\dd x\,\dd y)$\\
  $\mathcal C_0$, $\mathcal C_{j,\pm}$ & 中心胞元，第 $j$ 个左/右周期胞元（$j\geq0$）\\
  $\Gamma_\pm$, $x_\pm$, $\nu_\pm$ & 中心与半波导的人工界面 $\{x_\pm\}\times I_d$，
    $x_\pm=\pm L_0/2$，指向半波导的法向 $\pm e_x$\\
  $\Omega_0$, $\Omega_{j,\pm}$ & 中心与周期胞元中的介质柱\\
  $\mathcal G_\pm(k)$ & 半波导有限能量解在 $\Gamma_\pm$ 上的 Cauchy 数据空间\\
  $\psi_m$, $\beta_m$, $\gamma_m$, $K$ & 端口 Fourier 基，横向与纵向波数，通道数
    $K=2M+1$\\
  $(a,b)$ & 锚定在某个端口的右行、左行 Rayleigh 振幅\\
  $\tau$, $\omega$, $\eta$ & 双层势密度、单层势密度、$\eta=(\tau;-\omega)$\\
  $S,D,D',T$ & 单层、双层、伴随双层、超奇异边界积分算子\\
  $(\mathsf E_{\mathrm L},\mathsf E_{\mathrm R})$ & 单胞传播矩阵对\\
  $\mathsf Z_\pm$, $\mathsf P_\pm$ & 衰减子空间的振幅基，胞元传播矩阵\\
  $\mathsf V_\pm$ & 离散 Cauchy 数据基\\
  $\xi$, $c_\pm$ & 中心离散未知量，半波导衰减坐标\\
  $\mathsf M(k)$ & 中心离散非线性特征值问题的矩阵函数\\
  $\mathsf A(k)$ & 端口约化矩阵，候选由 $\det\mathsf A(k)=0$ 定位\\
  $\widehat u$ & 由候选构造的原始场\\
  $\jump{\widehat u}_\Gamma$, $\jump{\partial_\nu\widehat u}_\Gamma$ &
    界面上的值跳跃与法向导数跳跃\\
  $\varepsilon_g$, $C_g$, $\ell_g$ & 第 $g$ 组界面的加权缺陷、迹常数与参考长度\\
  $U$, $Q$ & 场范数，$Q=\sum_g C_g\varepsilon_g$\\
\end{longtable}
